% !TEX root = ../main.tex

この部では、Leanで仕様を証明するために必要な道具を整理する。ここでの目的は、定理をたくさん覚えることではない。むしろ、ソフトウェア開発でよく出会う主張を、どの種類の証明道具で扱えるのかを理解することである。

たとえば、リファクタリング前後の処理が同じであることを示したい場合、多くは等式推論になる。式を書き換え、不要な処理を消し、定義を展開し、最終的に同じ形に到達する。このとき使うのが \texttt{rw}、\texttt{simp}、\texttt{calc} である。これらは、数学の証明というより、レビュー可能な変形ログを書く道具として理解できる。

また、\texttt{Option}、\texttt{Bool}、\texttt{Sum} のような分岐を含む型では、場合分けが必要になる。エラーケースを忘れないことは、実務上も重要である。Lean の \texttt{cases} は、分岐を網羅的に扱うための道具として機能する。通常のテストでは代表例だけを見ることが多いが、証明では型が持つすべての形を分解して確認する。

リストや木構造のように、長さや深さが任意のデータを扱う場合には、帰納法が必要になる。帰納法は難しい数学の技法に見えるが、ソフトウェア的には「空のケースで成り立つ」「一つ追加しても成り立つ」なら「任意の長さで成り立つ」と考える方法である。バッチ移行で件数が保存されること、全要素が条件を満たすこと、木構造変換でノード数が保存されることなどに使える。

さらに、関数外延性、Galois接続、Pullback のような少し抽象的な道具も扱う。関数外延性は wrapper や adapter の同値性確認に使える。Galois接続は、詳細な情報から公開情報へ落とすような抽象化を説明するために使える。Pullback は、複数のデータソースが整合する条件、つまり join や制約付き合成を考えるための言葉になる。

この部を通じて、読者は「どの実務課題にはどの証明道具が対応するか」を見分けられるようになる。後続の仕様証明やマイグレーション証明では、この部の道具を具体的なサンプルに適用していく。

\section*{この部で扱うこと}

\begin{itemize}
\item 等式推論と書き換え
\item \texttt{simp} による正規化
\item \texttt{calc} による段階的な証明
\item \texttt{cases} による場合分け
\item \texttt{induction} によるリスト・木構造の証明
\item 関数外延性
\item 前順序とGalois接続
\item Pullback と整合制約
\end{itemize}

\section*{この部で扱わないこと}

\begin{itemize}
\item 自動証明探索の高度なチューニング
\item SMTソルバとの連携
\item 大規模なプログラム検証
\item 高度な圏論的証明
\end{itemize}

\section*{次の部への接続}

次の部では、ここで学んだ証明道具を使って、自然言語の仕様書をLean上の型・関数・命題・証明へ変換する。特に、口座操作や料金計算のような小さな業務仕様を題材に、証明可能な仕様の切り出し方を学ぶ。
